\begin{statementbox}
	\textbf{\textcolor{CStatement}{条件}}
	\begin{description}[style=nextline, leftmargin=2.8em, labelsep=0.5em, itemsep=0.35em]
		\item[\textbf{\textcolor{CStatement}{条件一（四点不共线）：}}] 四点$A,B,C,D$不共线（即任意三点不共线）。
	\end{description}

	\textbf{\textcolor{CStatement}{结论}}
	\begin{description}[style=nextline, leftmargin=2.8em, labelsep=0.5em, itemsep=0.45em]
		\item[\textbf{\textcolor{CStatement}{结论一（对角互补）：}}]
			\textbf{\textcolor{CStatement}{直观描述：}}
			四边形一组对角和为$180^{\circ}$。
			\[
				\angle A + \angle C = 180^{\circ}.
			\]

		\item[\textbf{\textcolor{CStatement}{结论二（同侧等角）：}}]
			\textbf{\textcolor{CStatement}{直观描述：}}
			线段$AB$同侧的两个角相等。
			\[
				\angle ACB = \angle ADB.
			\]

		\item[\textbf{\textcolor{CStatement}{结论三（相交弦逆定理）：}}]
			\textbf{\textcolor{CStatement}{直观描述：}}
			两线段$AB,CD$相交于点$P$，且$P$在线段内部，满足$PA \cdot PB = PC \cdot PD$。
			\[
				PA \cdot PB = PC \cdot PD.
			\]

		\item[\textbf{\textcolor{CStatement}{结论四（割线逆定理）：}}]
			\textbf{\textcolor{CStatement}{直观描述：}}
			两线段$AB,CD$的延长线相交于点$P$，且满足$PA \cdot PB = PC \cdot PD$。
			\[
				PA \cdot PB = PC \cdot PD.
			\]

		\item[\textbf{\textcolor{CStatement}{常见变形（圆幂逆定理）：}}]
			\textbf{\textcolor{CStatement}{直观描述：}}
			点$P$在直线$AB$和$CD$上，无论内分或外分，恒有$PA \cdot PB = PC \cdot PD$时四点共圆。
			\[
				PA \cdot PB = PC \cdot PD.
			\]
	\end{description}

\end{statementbox}
